\chapter{Open problems and proposed experiments}\label{sec:open}

% Survey section (drafted 2026-10-15, updated 2026-11-12). Each problem is phrased as the
% measurement or proof that would settle it, per the repository's
% "name the obstruction" rule.

Table~\ref{tab:fm-open} lists each problem with the measurement or proof
that would settle it and the value that measurement has now. The
paragraphs after it give the experiment in full.

\begin{table}[tbp]
\centering\scriptsize
\begin{tabularx}{\linewidth}{@{}r>{\raggedright}X>{\raggedright}p{0.33\linewidth}>{\raggedright\arraybackslash}p{0.13\linewidth}@{}}
\toprule
& settled by & current value (standing) & where \\
\midrule
1 & $d_{\rm ff}$, $d_F$ and solving degree of Weil-descended $\Sem{m+1}$, $m\ge3$, four or more $n$ per $(\ell,m)$ & $d_F\ge5$ at $n=33$--$45$ (certified, independently checked); $\ge4$ at $n=17$--$24$ (census, unverified) & Ch.~\ref{ch:falldegree}, \ref{ch:hardness} \\
2 & per-arm cold exponents, crossover only where fits separate across the band & no arm below matched rho; validated gaps $3.43$--$11.51\times$ (Table~A), $5.47$--$13.46\times$ prime and $10.62$--$38.21\times$ binary (Table~B) & \art{docs/ic/leaderboard.json} \\
3 & measured trials per relation against $m!\,r/N^{m-1}$ per base family & ECC2K-130 fit $C=0.27$ at $n=61$, $1.2$--$1.5$ at $n=71,73$, model $0.25$ (unvalidated) & \S\ref{sec:corrections}, item~3 \\
4 & a sub-$N^2$ non-generic membership test for $\Sem4$ over $\FF_p$, or a counting argument against one & none known; generic $k$-sum exponent $0.521$; generic single-target floor $1.5\times$ rho; sibling walk with \texttt{aut=3} $0.875$--$0.992\times$ rho-cubic at $w=0$ only (second points priced 0), $2.663$--$3.000\times$ at $w=1$ (validated GAE comparison) & Ch.~\ref{ch:distance}; Prop.~\ref{prop:collisionfloor}; Table~\ref{tab:distance-walk} \\
5 & IC advice--time curve $(S,T)$ against preprocessing rho at equal $S$ & IC $53$--$317\times$ slower online than Bernstein--Lange rho with an equal or smaller table & \S\ref{sec:ladder} \\
6 & $k=4$ replication, $m=14$--$17$, eight targets per arm, measured negation rho per rung & $5.2\times$ rho at $m=15$; crossover $m^*\approx17.6$ ($16.3$--$19.2$) extrapolated; E32 $1.5\cdot10^3\times$ at $2^{42.4}$ (validated) & Ch.~\ref{ch:extension}, \ref{ch:glvprogramme} \\
7 & the shared-rank arm's uncounted work, counted & three categories uncounted; the arm is a lower bound & \S\ref{sec:frontier} \\
8 & genus-3 GHS IC and negation rho at $q=2^{21}$--$2^{24}$, one cofactor, priced across the band & local exponent $1.356$ at $q=2^{22}$ (target ${<}4/3$); below the strongest rho at $q=2^{16}$ only under the registered price; $q=2^{23}$ exceeds the 4.5\,GB cap (no validation file) & \S\ref{sec:ladder2}, Ch.~\ref{sec:transfer} \\
9 & N4~v2 on an isolated host with a replayable rho baseline & exponent $0.983$; $0.005\times$ rho near $q=2^{21}$; $0.0021\times$ a rho model at $k=25$; factor-base constant $1.60$ costs $0.655$--$0.685$ of $2.00$ at $k=21$, $22$ (all L0, no validation file) & Ch.~\ref{sec:transfer}; \S\ref{sec:ladder4} \\
10 & a published index calculus (FGHR base, msolve, WDSat) as reference arm in ecbench & not run: no ratio against published IC at matched $(q,n,m)$ & \S\ref{sec:ladder3} \\
11 & sieve and Nagao arms with F4, descent and linear-algebra additions emitted, at $p\approx1300$--$2000$ on the isolated lab & sieve $0.0065\times$ matched negation rho at $p=1511$ and $0.0141\times$ pooled (both validated); sieve additions emitted at $0.816$ of the proxy (local); Nagao with Hilbert stop $1.003$, $0.837$, $0.695$ in U2 at $p=1301$, $1451$, $1601$ (local, lower bound) & \S\S\ref{sec:jac-jv},\ref{sec:ext-jv-nov} \\
12 & a proof that $d_F$ grows with $n$ or $k$ & no growing bound; best certified value 5 & Ch.~\ref{ch:falldegree} \\
13 & certificates for about 20 subspaces per $n$, $n=33$--$51$; guarded chain at $B=4$--$8$ & named instances only; guarded chain exact at $B=3$, $(4,5)$ & Ch.~\ref{ch:falldegree} \\
14 & Kummer-tower stage A at $t=1$--$14$ with closure certificates & session values $d_F=3,4,4,4,4,5,5,5,5,6$ at $t=1$--$10$ (not committed); F4 extrapolated to 1.7\,h, 17\,h, 7 days at $t=12$--$14$; the law $\lfloor(k+14)/4\rfloor$ quoted as measured to $k=15$ in a sibling repository; unchanged at characteristic 19--227 to $k=10$ (local, unverified) & \S\ref{sec:fd-pkm}; \S\ref{sec:open-leads} \\
\bottomrule
\end{tabularx}
\caption[]{Open problems, the measurement or proof that settles each, and
its current value. Values are copied from the sections named; standing is
as stated there, and an entry without a validation file or certificate is
a producer-reported value.}
\label{tab:fm-open}
\end{table}

\paragraph{1. The solving degree of Weil-descended summation systems.}
Every subexponential claim for prime-degree binary curves rests on a bound
for this degree (\S\ref{sec:theory}). No growing lower bound is proved;
Chapter~\ref{ch:falldegree} certifies $d_F\ge5$ at $n=33$--45 and the
census of Chapter~\ref{ch:hardness} gives $\ge4$ at $n=17$--24. The proof
half of the question is problem~12; this item is its measurement half.
\emph{Experiment:} measure first-fall, last-fall
and solving degree of $\Sem{m+1}$ restricted to $\ell$-dimensional
subspaces over a grid of $(n,\ell,m)$ with $m\ge3$, at least four values
of $n$ per $(\ell,m)$, against the semi-regular prediction, with
multi-graded regularity (one block per summand) as a second model.
\emph{Settles:} whether the degree grows with $n$ at fixed $(\ell,m)$; a
growth consistent with zero would revive the Petit--Quisquater and Semaev
heuristics, a positive slope would close them for this model.

\paragraph{2. Cold index calculus against rho with every operation
priced.} No measured arm on a prime-order or prime-degree binary curve
is below matched rho on the cold one-target boundary once native work is
priced on both sides (\S\ref{sec:frontier}).
\emph{Experiment:} the frontier campaign of \S\ref{sec:frontier} extended
to the largest sizes each arm can reach, with exponents fitted per arm and
the crossover reported only where the fitted exponents separate across the
whole pricing band. \emph{Settles:} whether any implemented arm has an
exponent below $1/2$ in $r$ on these families.

\paragraph{3. Relation yield is now the dominant stage.} Root finding made
the decomposition oracle cheap; what remains is the number of trials per
relation, governed by the factor-base density. \emph{Experiment:} for each
factor-base family, the measured trials-per-relation law against $r$ and
base size, compared with the counting prediction $m!\,r/N^{m-1}$; the
residual (e.g.\ the unexplained factor 5 in the ECC2K-130 extrapolation,
\S\ref{sec:koblitz}) is the target.

\paragraph{4. Root finding beyond two summands on prime curves.} The
two-summand oracle is prior art in mechanism (McGuire--Mueller) and costs
$\Theta(N)$ per trial; three summands cost a pair table of $N^2$ entries.
A generic table of signed 4-sums already moved the generic exponent to
$0.521$ (Chapter~\ref{ch:distance}), and any membership structure built
from group operations stays inside the $\sqrt r$ floor
(\S\ref{sec:gf-shoup}); a single-target generic arm with unknown base
logarithms costs at least $1.5\times$ rho
(Proposition~\ref{prop:collisionfloor}). \emph{Open:} whether a
sub-$N^2$ \emph{non-generic} (coordinate-algebraic) membership test for
$\Sem4$ over a prime-field base exists whose cost escapes that floor, for
example one built from the $D_m$-invariant coordinates of 2-torsion
curves. A counting argument against one would settle the question as well
as a construction.

\paragraph{5. The preprocessing regime.} Index calculus with a retained
table beats rho without one, and loses by $53$--$317\times$ to
Bernstein--Lange rho with an equal or smaller table (\S\ref{sec:ladder}).
\emph{Experiment:} measure the advice--time curve $(S,T)$ of the
compact-orbit arm across sizes against preprocessing rho's
$S\,T^2\approx N$; an index-calculus online result counts only if it lies
below that curve at equal $S$.

\paragraph{6. Symmetry at $k=4$, end to end.} The $D_4$ reduction and its
four-unknown weighted form moved a \emph{modelled} handover over
$\FF_{p^4}$ down from about $2^{150}$ (\S\S\ref{sec:ladder},\ref{sec:ladder2}).
Those single-point models are superseded: the E32 ladder, validated,
measured $1.5\cdot10^3\times$ its strongest matched rho at $n=2^{42.4}$,
falling as $n^{-0.24}$, with a fitted (not measured) crossover at
$2^{85.5}$ $[2^{79.2},2^{94.3}]$ (Chapter~\ref{ch:glvprogramme}). In
characteristic two a Frobenius-reduced form brought the gap to $4.3\times$
a measured rho at $n\approx2^{56}$ on one target ($6.1\times$ against the
rho model, \S\ref{sec:ladder2}). A red-team review
(Chapter~\ref{ch:extension}) found the extrapolated crossover near
$n\approx2^{68}$ unsupported, not refuted: the only clean point, $m=15$,
is $5.2\times$ rho after corrections, and the crossover extrapolates to
$m^*\approx17.6$ ($16.3$--$19.2$), at or above the edge of the data.
\emph{Experiment:} the frozen replication plan (eight targets per arm,
$m=14$--$17$, about 30\,h), with a measured negation-map rho at every
rung, cold and priced across the band; a measured crossover would be an
\cls{advance} on a weak family, and its absence would locate the term the
model misses.

\paragraph{7. Pricing and calibration infrastructure.} Rankings now
require every arm's native work to be counted (AGENTS.md \S5a).
\emph{Work:} count the three uncounted categories of the shared-rank arm;
calibrate every ladder size on one fixed host; register the $n=83$ dense
root index in the harness. (The cross-host F4 counter drift is closed:
since \#2043 every host takes the eager path, and sessions record the
version of the F4 accounting; Chapter~\ref{ch:audit}, item~4.)

\paragraph{8. Genus three after cubic GHS, to the crossover.} The one
family whose fitted exponent (about $q^{1.4}$) lies below rho's $q^{3/2}$ was
$1.15\times$ a measured negation-map rho at $q=2^{20}$ before rounds three
and four brought the whole solve to about $0.44\times$ the second round's
at $q=2^{22}$ (derived, in the inversion-charging unit;
\S\S\ref{sec:ladder},\ref{sec:ladder2}),
and is now below the strongest rho at $q=2^{16}$ under the registered
pricing, unresolved at the other sizes to $2^{20}$, and above it under a
heavier sensitivity pricing, with no measurable pooled trend
(\S\ref{sec:transfer}). What would settle it is the same comparison at
$q=2^{21}$--$2^{24}$ on instances of one cofactor, with the subgroup
multiplication's price measured rather than chosen;
its local exponent is down to $1.356$, still above $4/3$ because the linear
algebra grows faster than $q^{4/3}$, and $q=2^{23}$ needs more than the
4.5\,GB cap.
\emph{Experiment:} measure negation-map rho on the same instances in the
latest unit (inversions and mod-$\ell$ multiplications priced), and run
both past $q=2^{22}$, in the harness and across the pricing band. A measured
size crossover there would be an \cls{advance} against rho on a narrow and
undeployed family; the programme's one validated below-rho measurement
(the Joux--Vitse sieve, problem~11) is accounting at a single size.

\paragraph{9. Plane quartics: Diem's $\tilde O(q)$.} Genus-3 GHS covers whose
Jacobian is that of a smooth plane quartic admit Diem's index
calculus~\cite{diem2006,dt2008} at $\tilde O(q)$, against $\tilde O(q^{4/3})$ for
the hyperelliptic route and $q^{3/2}$ for rho on $E(\FF_{q^3})$. The
repository builds such quartics directly ($m=2$ GHS, toy index calculus at
$q\le11$) and through Smith's $(2,2,2)$-isogeny from the hyperelliptic case
(18.6\% of random curves); a four-workstream programme was scoped on 18
October (\art{research/plane_quartic_program_20261018/PROGRAM.md}).
\emph{Experiment:} the hyperelliptic ladder's instances, carried to plane
quartics where Smith's conditions pass, with Diem's relation search and the
same linear algebra, against the same matched rho. The 23 October literature
round ranks this branch as its one exponent candidate (LEAD-15,
\art{research/pdp_literature_leads_20261022/}).
\emph{Since then:} the arithmetic has been rebuilt with canonical class keys,
which resolve the degree-3 merges (\S\ref{sec:corrections}). On the direct
quartic route, index calculus is below a matched plain rho from $q=4051$,
with fitted exponent $1.083$ against $1.485$, and N4~v2 extends this
against a negation-map rho over $q=2^8$--$2^{22}$, at exponent $0.983$,
and $0.005\times$ rho at $q\approx2^{21}$ (L0, no validation file;
\S\ref{sec:transfer}). \emph{Still open:} a validation
of N4~v2 on an isolated host with a replayable baseline; index calculus on
the quartics reached through Smith's isogeny; and a record-scale run (for
example $q\approx2^{40}$, $\ell\approx2^{120}$) where rho is out of reach
and the published Diem--Thomé timings are the comparator.

\paragraph{10. A published index calculus as the reference arm.} Every
ratio in this treatise is against rho or the programme's own previous arm;
none is against the state of the art in index calculus at the same
$(q,n,m)$ (\S\ref{sec:ladder3}). \emph{Work:} fix, per curve and size, the
standard base $\{P:x(P)\in\FF_q\}$, the FGHR-symmetrised summation
polynomial (2-torsion and, on Koblitz curves, Frobenius invariants), msolve
as the solver and WDSat for binary decompositions, and run it in ecbench
on the existing ladders (binary $n=13$--$37$, Koblitz $n=41$--$61$, $k=4$
extension fields). \emph{Settles:} whether any engineering result of
\S\S\ref{sec:rootfinding}--\ref{sec:transfer} improves on published index
calculus, as opposed to on this programme's own baseline. Stage ratios
quoted against the in-house F4 say nothing about published Gr\"obner
engines until this is done.

\paragraph{11. The Joux--Vitse cover, to the crossover.} On 24 October the
plain Nagao arm of the cover of $E(\FF_{p^6})$ was $1.6$--$2.0\times$ a
pooled rho constant at $p=1009$ (a lower bound), closing as about
$p^{-2}$, with the line meeting rho near $p\approx1330$ by extrapolation
(\S\ref{sec:ladder3}). The November studies ran rho on the same instances
(\S\ref{sec:jac-jv}), and the extrapolation does not hold. Against a
matched negation rho the plain arm is above rho in group additions at
$p=1301$ ($1.405$) and $1451$ ($1.160$) and unresolved at $1601$
($0.870$ $[0.697,1.183]$), and above it in $\FF_p$ multiplications at all
three (\art{research/jv_cover_negrho_rescore_20261101/}). A factor-base
line search makes it $1.71\times$ cheaper, below rho in additions
($0.791$--$0.548$) but above in multiplications ($1.900$--$1.318$), so it
is unranked across the band; its modelled crossovers are $p^*\approx1185$
in additions and $1799$ in multiplications, with no validation file
(\art{research/jv_cover_fbline_20261102/}). On 12 November the Hilbert
early stop of F4 halved the line-search arm's counted cost at equal
relations, and the arm moved to $1.003$, $0.837$ and $0.695\times$ rho in
multiplications at $p=1301$, $1451$ and $1601$, below across the band at
the two larger primes (local, no validation file;
\S\ref{sec:ext-jv-nov}). The \emph{sieve} arm is
$0.0065\times$ a matched negation rho at $p=1511$ in group additions
($155\times$ $[125,184]$ below it; $64\times$ in $\FF_p$
multiplications; at most $0.12\times$ across the band), on the weak class
only, validated as accounting (\texttt{VAL-20261103-c31c76}); it has no
\texttt{pcc} claim record, so its statement is a validation-gate result
and otherwise an unverified assertion. Five questions remain:
\begin{itemize}
\item whether the sieve arm's order survives once its F4, descent and
  linear-algebra additions are emitted and priced, rather than bounded
  (emitted locally on 12 November at $0.816$ of the proxy, every cell
  still below rho; not yet on the isolated lab);
\item whether the Nagao row with line search and Hilbert stop validates,
  in both units, with its $\FF_p$ additions emitted;
\item whether the route reaches the weak class from a typical curve: the
  implemented isogeny walk reaches a weak curve from 3 of 60 random
  admitted starts at $p=1009$ and none at $1511$ and $1823$, and odd
  isogenies with irrational kernels are not implemented;
\item why the two arms differ by about $30\times$ at $p=1009$;
\item whether both results hold on the isolated lab rather than a local
  run.
\end{itemize}
\emph{Experiment:} emit those counters on both arms, and run both arms and
rho at $p\approx1300$--$2000$ on the isolated lab, across the band. On a
weak, undeployed class this would turn a validated accounting result into a
priced result.

\paragraph{12. A growing lower bound on the last fall degree.} Chapter~\ref{ch:falldegree}
proves the first fall degree constant and certifies $d_F\ge5$ on named
binary instances, but no presentation has a proved lower bound that grows
with $n$ or $k$. \emph{Proof route:} construct degree-$c$ pseudo-solutions,
i.e.\ linear functionals on $B_{\le c}$ that kill $W_c(F)$ and not 1, for a
family of instances, generalising the closure certificates of
Proposition~\ref{prop:cert}. The regular closure pattern observed (exactly
$4n$ new cubic dimensions per round, then a stall) and the
iterated-system lower-bound technique of~\cite{iterated2023} are the
starting points. \emph{Settles:} the open problem of
HKY~\cite{hky2015} and Kousidis--Wiemers~\cite{kousidis2019}; a bound
$\omega(\sqrt{n\ln n})$ would invalidate the subexponential conclusion of
Semaev's analysis for these systems.

\paragraph{13. Certificates at scale.} \emph{Experiment:} certify
$1\notin W_4$ (and, where it holds, $1\notin W_5$) on about ten random and
ten polynomial subspaces per $n$ for $n=33$--$51$, publish the
certificates, and extend the guarded chain to $B=4$--$8$ by computation.
\emph{Settles:} whether the step $4\to5$ is a population property, and
where $5\to6$ occurs; the exact $d_F$ of the guarded chain beyond $B=3$.

\paragraph{14. The prime-field Kummer tower, re-measured.} The session
values $d_F=3,4,4,4,4,5,5,5,5,6$ at $t=1..10$ were never committed.
\emph{Experiment:} regenerate stage A at $t=1..14$ with closure
certificates and score the registered prediction
$d_F=4+\lfloor(t-2)/4\rfloor$; extrapolated F4 times are about 1.7\,h, 17\,h
and 7 days at $t=12$--$14$ (\art{research/prime_ic_program_20261029/}).
\emph{Settles:} whether the proved constant first fall degree 4 sits beside
a last fall degree that keeps growing in the setting PKM introduced.
The registered prediction equals $\lfloor(t+14)/4\rfloor$, the law that
the 12 November survey quotes as measured to $k=15$ in the sibling
\texttt{ml-cryptanalysis} repository, with a written proof of unbounded
growth at slope $\ge1/14$ (quoted, not checked here;
\art{research/leads_20261112/README.md}). The chase of
\S\ref{sec:open-leads} found the law unchanged with one equation at
characteristic 19--227 to $k=10$, and one degree lower with a second
equation.

\section{Leads opened, run, closed or killed, 11--12 November}\label{sec:open-leads}

Three lead surveys and one chase ran on 11 and 12 November at the owner's
request. Each lead was preregistered with a falsifier before its test; the
skeptic's reading replaces the producer's where they disagree. Every
outcome is \texttt{[pcc:unverified]}: local runs, no pcc claim, no
validation file. The novelty checks compared each lead with the 173 lead
ids of the 10 November sweep and the 176 ids of the crypto-autoresearcher
lead atlas; the atlas bodies were vault-encrypted and the
\texttt{research-kb} server unreachable, so ``new'' means new with respect
to those ids. The atlas itself was last reviewed on 24 October and does
not yet record these dispositions.

\paragraph{The 11 November survey.} Three idea-generator lanes proposed 13
leads (\art{research/leads_20261111/README.md}): lane~A on the last fall
degree and the selector proof system, lane~B on binary-field engineering,
lane~C on prime and extension fields and covers. All 13 and one follow-up
were tested on 12 November (Table~\ref{tab:open-leads1111}). Three bear on
rows: the Nagao Hilbert stop held, the sieve emission was partial
(accounting), and the sieve Hilbert stop was killed
(\S\ref{sec:ext-jv-nov}). Lane~A's results bear on problems~12 and 13, not
on any row.

\begin{table}[tbp]
\centering\scriptsize
\begin{tabularx}{\linewidth}{@{}>{\raggedright}p{0.30\linewidth}>{\raggedright}p{0.13\linewidth}>{\raggedright\arraybackslash}X@{}}
\toprule
Lead & Outcome & Decisive measurement (rule) \\
\midrule
\texttt{LEAD-C-jv-nagao-hilbert-at-scale} & held & $S_B/S_A=0.5276$--$0.5280$ at equal digests; U2 $1.003$, $0.837$, $0.695$ of rho (falsifiers F1--F4 pass) \\
\texttt{LEAD-C-jv-sieve-emit-price} & partial & emitted/proxy $0.816$ (ratio of means; per seed $0.44$--$0.93$); kill line $1.20$ not reached \\
\texttt{LEAD-C-jv-sieve-hilbert-descent} & killed & descent native F4 $1.139\times$ (kill line $1.10$); band order flips at $w=1$ \\
\texttt{LEAD-C-quartic-n4v2-k26} & inconclusive & $k=25$ reached, $k=26$ over the 3\,GB limit; exponent lower bound $1.1000$--$1.1034$ at the $1.10$ kill line (skeptic) \\
\addlinespace
\texttt{LEAD-A-selector-eta-kernel} & falsified & nullity exceeds the trivial kernel at feasible cells: $1143$--$1145$ vs $232$ ($n=23$, $e=4$) \\
\texttt{LEAD-A-selector-counting-law} & held & model A 4 of 5 cells, model B 2 of 8 (preregistered grid) \\
\texttt{LEAD-A-fall-effective-variable-law} & falsified & $n=35$ ($v=34$) and $n=36$ ($v=35$) fall at $D=4$, 2 of 2 each, against ``$1\in W_4$ iff $v\le33$'' \\
\texttt{LEAD-A-closure-dimension-laws} & held & $\dim(W_3\cap B_{\le3})=n(5n+1)/2$ at $n=23$--$33$; $\dim(W_4\cap B_{\le3})=3n(n+2)$ at $n=35$--$43$ \\
\texttt{LEAD-A-bilinear-lift-size-degree} & partial & $d_F=3$ at $n=11$, 13, 15; the random negative control also drops $4\to3$ \\
\texttt{FOLLOWUP-selector-extra-syzygies} & held & extra $c$-only syzygies only at $e=e_{\min}$; excess 0 at $e_{\min}-1$ on 18 of 18 \\
\addlinespace
\texttt{LEAD-ICB-tableB-binary-\allowbreak known-log-claw} & prediction failed & $0.88$--$1.44\times$ rho-neg at f2m15--27 against $[1.4,2.4]$; a BSGS-class boundary result, outside the IC ratchet (owner, 8 November) \\
\texttt{LEAD-ICB-d4bin-chain-variable-sweep} & closed & 12 chains, not 16; the alternative gives 28 chains at $2.34\times$ per trial \\
\texttt{LEAD-ICB-d4bin-exponent-coset} & closed & difference rank 4 on 300 of 300: no index-2 coset \\
\texttt{LEAD-ICB-d4bin-constant-\allowbreak quadratic-elimination} & falsified as written & quadratic part rank 2, not 3, on 300 of 300 \\
\bottomrule
\end{tabularx}
\caption[]{Dispositions of the 11 November leads, tested 12 November.
Outcome: the skeptic's reading where it differs from the producer's
(\texttt{known-log-claw}, \texttt{quartic-n4v2-k26}). All local and
unverified. Sources: \art{research/leads_20261111/README.md} and
\art{research/leads_20261111/runs/}; \art{research/jv_nagao_hilbert_20261112/};
\art{research/jv_sieve_hilbert_descent_20261112/};
\art{research/lfd_lower_bound_20261111/out/}.}
\label{tab:open-leads1111}
\end{table}

\paragraph{The 12 November survey and its chase.} A second survey, by
field regime, proposed 15 techniques, dropped P4 (a misreading of the
isogeny-tower cell) and folded M2 into P2 (\art{research/leads_20261112/README.md}). It
first fixed a bar: with $m$ summands in a coset of $\mu_{2^k}$, a Gröbner
route on the chain presentation beats brute-force decomposition only if
$\mathrm{LFD}/k<0.184$ at $m=3$ with sparse linear algebra ($0.083$ at
$m=2$). A chase ran the top five under a frozen protocol
(\art{research/leads_chase_20261112/RESULTS.md}):
\begin{itemize}
\item \emph{J2} (percolation factor-base size) gave a held-out measured
  improvement on the N4~v2 row: $0.655$--$0.685$ of the charged cost at
  $c=1.60$, $k=21$ and $22$, falsified at $k=16$ (\S\ref{sec:ladder4});
\item \emph{S1} (characteristic and generator crossing) left the number of
  equations (Dm) as the only surviving hypothesis: a second equation over
  12289 puts $1$ in $V_5$ at $k=10$, one degree below the law's 6, while characteristic $19$--$227$
  with one equation leaves the law on every scored cell; the record's
  Fermat-prime cells were void, since those fields admit no unit instance;
\item \emph{P1} (penultimate quotient law) reproduced the committed
  $q/2^k$ exactly on 72 of 72 cells at three fresh primes, which confirms
  instance independence; its decisive stage, $D=6$ at $k=15$, needs a GPU;
\item \emph{M1} (three-summand arity) closes at the measured depths from
  existing records: $\mathrm{LFD}/k=1.4$ at $k=5$ and $\ge1.33$ at $k=6$, with
  16 and 26 times more columns than enumeration;
\item \emph{J1} (the quartic solver on generic Jacobians) is blocked: the lab
  has no genus-3 point counter, and $\ell<2^{64}$ caps instances near
  $q\approx2^{21}$.
\end{itemize}
P2, P3, P5, J3, J4, S2, S3 and K1--K3 were proposed and not run. K1, a
Frobenius fold on the twisted line of $E/\FF_{16}$, predicts the $k=4$
crossover of problem~6 moving from $m^*\approx17.6$ to about $15.5$
$[14.2,17.1]$; it needs a matched-rho arm first. The survey also proposed
atlas verdict changes for cells 0.1, 0.2, E3, D4, C5 and A4; they are
proposals, not verdicts.

\paragraph{Curve25519, wave five.} A third survey proposed 13 leads for
Curve25519 index calculus and ran none
(\art{research/curve25519_leads_wave5_20261112/README.md}). It checked one
law against committed cells: the Gröbner solving degree of a decomposition
system on a subgroup base $\mu_d$ equals $d+1$ at every prime $d$ (20 of 20
cells), and $e_{\max}+2$ to $3$ ($m=2$) or $e_{\max}+3$ to $5$ ($m=3$) at
composite $d$, with $e_{\max}$ the largest prime factor; eight
preregistered $k=22$, $m=3$ values matched. For the real curve,
$p-1=4\cdot3\cdot65147\cdot C$ with $C$ a 236-bit probable prime, so every
$\mu_d$ base with $d\nmid12$ has solving degree at least $65{,}148$ on that
route. That is a security statement for the subgroup-base route, not an
attack. The two cheapest improvement leads predict $-2\%$ (a final check
from the stored ladder) and $0.83\times$ at $k=20$ (storing set-up points),
neither changing an exponent; after all four pipeline levers the source
puts the $C=2$ track at about $1.7$--$1.8\times$ rho-neg at $k\ge28$.
